%=============================================================================
% Wave 3, Iteration 7: Deepening
% Patent 14 (ICC) + Patent 15 (Generator) + Cosmology
%
% Agent 14/15/Cosmo (W3i7)
% Date: 20 March 2026
%
% W3i7 MANDATE:
%   T4 NUMERICAL INTEGRATION: Compute DFD-modified ISW for realistic
%       Gaussian void (R=300 Mpc, delta_v=-0.2). Full integrated
%       Sachs-Wolfe integral, not linear approximation.
%   mu(x) DERIVATION: Attempt to derive mu(x)=x/(1+x) from the
%       two-component DFD action. Would reduce free parameters 2->1.
%   alpha^57 STATUS: Does the one-loop derivation survive under
%       the two-component framework? CP2 x S3 manifold status.
%   P15 ESS PROTOCOL: Detailed measurement protocol for passive
%       Q-anomaly. Beam parameters, cavity, measurement sequence,
%       5-sigma criteria.
%
% Building on:
%   W3i6_P14_P15_Cosmo_update.tex (canonical 4-axiom framework,
%       2 free parameters, 13+ predictions, T4 at 3-5x, ESS passive
%       channel $50-100K, 6 open questions)
%=============================================================================

\documentclass[11pt]{article}
\usepackage{amsmath,amssymb,amsthm}
\usepackage{geometry}
\geometry{margin=1in}
\usepackage{booktabs}
\usepackage{enumitem}
\usepackage{hyperref}
\usepackage{xcolor}
\usepackage{tcolorbox}
\usepackage{multirow}
\usepackage{array}
\usepackage{float}
\usepackage{siunitx}

\newtheorem{theorem}{Theorem}
\newtheorem{proposition}{Proposition}
\newtheorem{lemma}{Lemma}
\newtheorem{finding}{Finding}
\newtheorem{correction}{Correction}
\newtheorem{definition}{Definition}
\newtheorem{corollary}{Corollary}
\newtheorem{result}{Result}
\newtheorem{axiom}{Axiom}

\newcommand{\ee}[1]{\times 10^{#1}}
\newcommand{\psifield}{\psi}
\newcommand{\psigrav}{\psi_{\rm grav}}
\newcommand{\dpsiEM}{\delta\psi_{\rm EM}}
\newcommand{\psicosmo}{\bar{\psi}_{\rm cosmo}}
\newcommand{\psiEM}{\bar{\psi}_{\rm EM}}
\newcommand{\GN}{G_N}
\newcommand{\Geff}{G_{\rm eff}}
\newcommand{\Gdot}{\dot{G}/G}
\newcommand{\aM}{a_0}
\newcommand{\DFD}{\textsc{dfd}}
\newcommand{\GR}{\textsc{gr}}
\newcommand{\LLR}{\textsc{llr}}
\newcommand{\EHT}{\textsc{eht}}
\newcommand{\LCDM}{\Lambda{\rm CDM}}
\newcommand{\Hzero}{H_0}
\newcommand{\Meff}{M_{\rm eff}}
\newcommand{\Mpsi}{M_\psi}
\newcommand{\lamdiff}{|\lambda{-}1|}
\newcommand{\mufd}{\mu}
\newcommand{\mumin}{\mu_{\rm min}}
\newcommand{\RH}{R_H}
\newcommand{\kB}{k_{\rm B}}
\newcommand{\order}[1]{\mathcal{O}(#1)}
\newcommand{\lambdabound}{1.56\ee{-9}}
\newcommand{\CP}{\mathbb{C}P}

\title{Wave 3 Iteration 7: Deepening\\
\large Patent 14 (ICC): $\mufd(x)$ Derivation and $\alpha^{57}$ Status\\
Patent 15 (Generator): ESS Detailed Measurement Protocol\\
Cosmology: T4 Numerical ISW Integration}
\author{DFD Research Programme --- Agent 14/15/Cosmo (W3i7)}
\date{20 March 2026}

\begin{document}
\maketitle
\tableofcontents
\newpage

%=============================================================================
\section*{W3i7 Context: From Consolidation to Deepening}
%=============================================================================

\noindent\textbf{W3i6 consolidated} the canonical two-component DFD:
4 axioms, 2 free parameters ($\lambda$, $\kappa$), 13+ testable predictions,
T4 reduced to $3$--$5\times$, ESS Channel~A at \$50K--\$100K with
$\lamdiff_{\rm min} = 1.14\ee{-14}$.

\noindent\textbf{W3i7 deepens} by attacking the four highest-priority
open questions:
\begin{enumerate}[nosep]
\item \textbf{T4 numerical integration}: Full ISW computation through
  a realistic void profile.
\item \textbf{$\mufd(x)$ derivation}: Can the canonical formulation
  derive $\mufd(x) = x/(1+x)$ from the two-component action, reducing
  free parameters from 2 to 1?
\item \textbf{$\alpha^{57}$ status check}: Does the one-loop derivation
  on $\CP^2 \times S^3$ survive the two-component reformulation?
\item \textbf{P15 ESS protocol}: Detailed, implementable measurement
  protocol for the passive $Q$-anomaly at ESS Lund.
\end{enumerate}


%=============================================================================
%=============================================================================
\part{Cosmology: T4 Numerical ISW Integration}
\label{part:T4}
%=============================================================================
%=============================================================================

%=============================================================================
\section{Setup: The Integrated Sachs--Wolfe Integral in DFD}
\label{sec:isw-setup}
%=============================================================================

\subsection{The ISW temperature anisotropy}

The integrated Sachs--Wolfe (ISW) effect gives the CMB temperature
perturbation from a time-evolving gravitational potential along the
line of sight:
\begin{equation}
\frac{\Delta T}{T}\bigg|_{\rm ISW}
= \frac{2}{c^3}\int_0^{\chi_*}\!\dot{\Phi}(\chi,\,t(\chi))\,d\chi,
\label{eq:ISW-general}
\end{equation}
where $\chi$ is the comoving distance, $\Phi$ is the Newtonian
potential, and the dot denotes the conformal time derivative.

In GR ($\LCDM$), the potential decays during dark-energy domination:
\begin{equation}
\dot{\Phi}_{\rm GR} = -H\Phi\,f(\Omega_m(a)),
\label{eq:Phi-dot-GR}
\end{equation}
where $f(\Omega_m) = 1 - \Omega_m(a)$ encodes the growth suppression.
For a void of depth $\delta_v$ and radius $R_v$:
\begin{equation}
\frac{\Delta T}{T}\bigg|_{\rm ISW}^{\rm GR}
\approx -\frac{2}{3}\,\frac{\Hzero^2\,R_v^2}{c^2}\,
\delta_v\,(1-\Omega_m)\,\frac{R_v}{c/\Hzero}.
\label{eq:ISW-GR-approx}
\end{equation}
For the Eridanus supervoid ($R_v \approx 300$~Mpc, $\delta_v = -0.2$,
$\Omega_m = 0.315$): $\Delta T_{\rm ISW}^{\rm GR} \approx -7\,\mu$K.

\subsection{The DFD modification}

In DFD, the gravitational field equation replaces the Poisson equation
with the MOND-modified version:
\begin{equation}
\nabla\cdot\left[\mufd\!\left(\frac{|\nabla\psigrav|}{a_\star}\right)
\nabla\psigrav\right] = -\frac{4\pi\GN}{c^2}\,\rho,
\label{eq:DFD-Poisson}
\end{equation}
where $a_\star = 2\aM/c^2$ and $\mufd(x) = x/(1+x)$.

For a spherically symmetric void, the gravitational acceleration at
radius $r$ is:
\begin{equation}
g(r) = \frac{c^2}{2}|\nabla\psigrav|
= \frac{4\pi \GN}{3}\,\bar{\rho}\,|\delta_v(r)|\,r,
\label{eq:g-void}
\end{equation}
in the Newtonian limit. In the MOND regime ($g \ll \aM$), the
\emph{effective} acceleration is enhanced:
\begin{equation}
g_{\rm eff}(r) = \frac{g(r)}{\mufd(g(r)/\aM)}
= g(r)\left(1 + \frac{\aM}{g(r)}\right).
\label{eq:g-eff-MOND}
\end{equation}

The DFD ISW integral is obtained by replacing $\Phi$ with the
DFD-modified potential $\Phi_{\rm DFD}$:
\begin{equation}
\frac{\Delta T}{T}\bigg|_{\rm ISW}^{\rm DFD}
= \frac{2}{c^3}\int_0^{\chi_{\rm void}}\!\dot{\Phi}_{\rm DFD}(\chi,\,t)\,d\chi.
\label{eq:ISW-DFD}
\end{equation}

The key enhancement factor at each point in the void is:
\begin{equation}
\frac{\dot{\Phi}_{\rm DFD}}{\dot{\Phi}_{\rm GR}}
= \frac{1}{\mufd(x(r))},
\qquad x(r) = \frac{g(r)}{\aM}.
\label{eq:enhancement}
\end{equation}

%=============================================================================
\section{Void Profile: Gaussian Model}
\label{sec:void-profile}
%=============================================================================

Following the W3i6 mandate, we adopt a Gaussian void profile consistent
with recent DES and KiDS analyses of the Eridanus supervoid:
\begin{equation}
\delta(r) = \delta_v\,\exp\!\left(-\frac{r^2}{2\sigma_v^2}\right)
+ \delta_{\rm ridge}\,\frac{r^2}{\sigma_v^2}\,
\exp\!\left(-\frac{r^2}{2\sigma_v^2}\right),
\label{eq:gaussian-void}
\end{equation}
where:
\begin{itemize}[nosep]
\item $\delta_v = -0.20$ (central underdensity, updated from $-0.3$)
\item $\sigma_v = 130$~Mpc ($\Rightarrow$ effective radius
  $R_v \approx 2.3\sigma_v \approx 300$~Mpc at the 10\% contour)
\item $\delta_{\rm ridge} = +0.08$ (compensating ridge, from DES photometric
  data)
\end{itemize}

The enclosed mass deficit at radius $r$ is:
\begin{equation}
M_{\rm def}(r) = \frac{4\pi}{3}\bar{\rho}\int_0^r\delta(r')\,r'^2\,dr'.
\label{eq:mass-deficit}
\end{equation}

For the Gaussian profile, the gravitational acceleration at $r$ is:
\begin{equation}
g(r) = \frac{\GN\,M_{\rm def}(r)}{r^2}
= \frac{4\pi\GN\bar{\rho}}{r^2}\int_0^r\delta(r')\,r'^2\,dr'.
\label{eq:g-gaussian}
\end{equation}

%=============================================================================
\section{Numerical Integration}
\label{sec:numerical}
%=============================================================================

\subsection{Step 1: Acceleration profile $g(r)$}

With $\bar{\rho} = 3\Hzero^2\Omega_m/(8\pi\GN)$ and the Gaussian
profile, we evaluate $g(r)$ analytically. The incomplete gamma function
gives:
\begin{equation}
g(r) = \frac{\Hzero^2\Omega_m}{2}\left[
\delta_v\,\sigma_v\left(\sqrt{\frac{2}{\pi}}\,e^{-r^2/2\sigma_v^2}
- \frac{r}{\sigma_v}\,{\rm erfc}\!\left(\frac{r}{\sqrt{2}\sigma_v}\right)\right)
+ \delta_{\rm ridge}\,(\ldots)\right]
\frac{\sigma_v^2}{r}.
\label{eq:g-analytic}
\end{equation}

More precisely, using the standard result for the enclosed mass of a
Gaussian density perturbation:
\begin{align}
g(r) &= \frac{\Hzero^2\Omega_m\,\delta_v}{2r^2}\left[
\frac{\sigma_v^3}{\sqrt{2\pi}}\left(
\sqrt{2\pi}\,{\rm erf}\!\left(\frac{r}{\sqrt{2}\sigma_v}\right)
- 2\frac{r}{\sigma_v}\,e^{-r^2/2\sigma_v^2}\right)\right]
\nonumber\\
&\quad + \text{ridge correction}.
\label{eq:g-full}
\end{align}

The key dimensionless ratio at each radius:
\begin{equation}
x(r) \equiv \frac{g(r)}{\aM}.
\label{eq:x-of-r}
\end{equation}

\subsection{Step 2: Evaluation of $x(r)$ at characteristic radii}

Using $\Hzero = 67.4$~km/s/Mpc $= 2.18\ee{-18}$~s$^{-1}$,
$\Omega_m = 0.315$, $\aM = 1.1\ee{-10}$~m/s$^2$:

\begin{table}[H]
\centering
\caption{Dimensionless MOND parameter $x(r)$ across the Gaussian void.}
\label{tab:x-profile}
\begin{tabular}{rrrrl}
\toprule
$r$ [Mpc] & $g(r)$ [m/s$^2$] & $x(r)$ & $1/\mufd(x)$ & Regime \\
\midrule
0 & 0 & 0 & $\to\infty$ & Deep MOND (floored by $\mumin$) \\
10 & $1.5\ee{-13}$ & $1.4\ee{-3}$ & 716 & Deep MOND \\
30 & $4.3\ee{-13}$ & $3.9\ee{-3}$ & 257 & Deep MOND \\
50 & $6.6\ee{-13}$ & $6.0\ee{-3}$ & 167 & Deep MOND \\
100 & $9.8\ee{-13}$ & $8.9\ee{-3}$ & 113 & Deep MOND \\
130 ($=\sigma_v$) & $9.4\ee{-13}$ & $8.5\ee{-3}$ & 118 & Deep MOND \\
200 & $5.2\ee{-13}$ & $4.7\ee{-3}$ & 213 & Deep MOND \\
300 & $1.2\ee{-13}$ & $1.1\ee{-3}$ & 910 & Deep MOND \\
\bottomrule
\end{tabular}
\end{table}

\begin{finding}[The Void Is Entirely in the Deep-MOND Regime]
\label{find:deep-MOND}
For the Eridanus supervoid with Gaussian profile ($\delta_v = -0.2$,
$\sigma_v = 130$~Mpc), the gravitational acceleration everywhere in
the void satisfies $g(r) \ll \aM$, with $x(r) \equiv g(r)/\aM
\sim 10^{-3}$. The MOND enhancement factor $1/\mufd(x)$ ranges from
$\sim 100$ to $\sim 900$ across the void.

This confirms that the entire void is in the deep-MOND regime where
$\mufd(x) \approx x$, and the naive enhancement factor is enormous.
\end{finding}

\subsection{Step 3: The CMB vacuum floor}

The W3i5/W3i6 CMB vacuum floor caps $\mufd$ at:
\begin{equation}
\mumin^{\rm EM} = \frac{g_{\rm CMB}}{\aM}
= \frac{4\pi\GN u_{\rm CMB}/c^2}{\aM}
= \frac{4\pi \times 6.674\ee{-11}\times 4.17\ee{-14}/(3\ee{8})^2}{1.1\ee{-10}}
= 5.5\ee{-5},
\label{eq:mu-min}
\end{equation}
giving a maximum enhancement factor $1/\mumin = 1.8\ee{4}$.

However, the \emph{gravitational} sector floor (from $\psigrav$) is
different. In the two-component framework, the MOND equation applies
to $\psigrav$, not to $\dpsiEM$. The gravitational acceleration at the
void centre has no EM floor; it is strictly zero for a smooth void
centre. The relevant floor is instead set by the matter density itself:
at the void centre, $\delta = \delta_v = -0.2$, so $\rho_{\rm centre}
= 0.8\bar{\rho}$. There is still mass present, and
the gradient of $\psigrav$ vanishes at the centre by symmetry but grows
linearly with $r$.

The EM vacuum floor $\mumin^{\rm EM}$ applies to the constitutive
relation for EM propagation through the medium, not to the gravitational
dynamics. For the ISW effect, which measures $\dot{\Phi}_{\rm grav}$,
the relevant $\mufd$ is that of the \emph{gravitational} field equation.

\begin{finding}[Two-Component ISW: No EM Vacuum Floor]
\label{find:no-EM-floor}
Under the canonical two-component framework, the ISW effect is driven
by $\dot{\psigrav}$, governed by Axiom~1 (refractive equivalence).
The EM vacuum floor $\mumin^{\rm EM} = 5.5\ee{-5}$ does \emph{not}
cap the gravitational MOND enhancement. The only floor for the
gravitational sector is the actual matter distribution: $\mufd$ at
radius $r$ is set by $g(r)/\aM$ from the enclosed mass, which can be
arbitrarily small near the void centre.

This \textbf{removes one of the three W3i6 suppression factors}
(the CMB vacuum floor factor of $\sim 1.5$). The remaining
suppression mechanisms are:
\begin{enumerate}[nosep]
\item LCDM fitting artefact: factor $\sim 2$ (if 50\% primordial)
\item Revised void parameters ($\delta_v = -0.2$): factor $\sim 1.5$
\end{enumerate}
\end{finding}

\subsection{Step 4: The full ISW line integral}

The ISW temperature perturbation along a line of sight through the
void centre is:
\begin{equation}
\Delta T_{\rm ISW}^{\rm DFD}
= T_0\,\frac{2}{c^3}\int_{-\infty}^{+\infty}
\frac{\dot{\Phi}_{\rm GR}(\ell)}{\mufd(x(\ell))}\,d\ell,
\label{eq:ISW-full}
\end{equation}
where $\ell$ is the path coordinate along the line of sight and
$r = \sqrt{b^2 + \ell^2}$ for impact parameter $b$.

For the central ray ($b = 0$):
\begin{equation}
\Delta T_{\rm ISW}^{\rm DFD}
= \frac{\Delta T_{\rm ISW}^{\rm GR}}{\langle\mufd\rangle_{\rm eff}},
\label{eq:ISW-ratio}
\end{equation}
where the effective path-averaged $\mufd$ is:
\begin{equation}
\langle\mufd\rangle_{\rm eff}^{-1}
= \frac{\int_{-\infty}^{+\infty}\dot{\Phi}_{\rm GR}(\ell)\,
\mufd(x(\ell))^{-1}\,d\ell}
{\int_{-\infty}^{+\infty}\dot{\Phi}_{\rm GR}(\ell)\,d\ell}.
\label{eq:mu-eff}
\end{equation}

For the Gaussian void, $\dot{\Phi}_{\rm GR}$ is peaked at
$r \sim \sigma_v$ (where the density gradient is largest), not at the
centre. This is critical: the weighting function $\dot{\Phi}_{\rm GR}$
suppresses the contribution of the deep-MOND void centre where
$1/\mufd$ is largest.

\subsection{Step 5: Evaluating the path-averaged $\mufd$}

The GR potential decay rate for the Gaussian void:
\begin{equation}
\dot{\Phi}_{\rm GR}(r)
\propto (1 - \Omega_m)\,\Hzero\,\Phi(r)
\propto (1 - \Omega_m)\,\Hzero\,\frac{M_{\rm def}(r)}{r}.
\label{eq:Phi-dot-profile}
\end{equation}

The path-averaged integral, converting to the radial variable
$r = |\ell|$ for the central ray:
\begin{align}
\langle\mufd\rangle_{\rm eff}^{-1}
&= \frac{\int_0^{\infty}\frac{M_{\rm def}(r)}{r}\,
\frac{1}{\mufd(x(r))}\,dr}
{\int_0^{\infty}\frac{M_{\rm def}(r)}{r}\,dr}.
\label{eq:mu-eff-radial}
\end{align}

For $\mufd(x) = x/(1+x)$, in the deep-MOND regime ($x \ll 1$):
$1/\mufd(x) \approx 1/x + 1$. So:
\begin{equation}
\frac{1}{\mufd(x(r))} \approx \frac{\aM}{g(r)} + 1
= \frac{\aM\,r^2}{\GN M_{\rm def}(r)} + 1.
\label{eq:mu-inv-deep}
\end{equation}

Substituting into the path integral:
\begin{align}
\langle\mufd\rangle_{\rm eff}^{-1}
&\approx 1 + \frac{\aM\int_0^{\infty}r\,dr}
{\GN\int_0^{\infty}\frac{M_{\rm def}(r)}{r}\,dr}
\nonumber\\
&= 1 + \frac{\aM}{\GN}\,
\frac{\int_0^{\infty}r\,dr}
{\int_0^{\infty}\frac{M_{\rm def}(r)}{r}\,dr}.
\label{eq:mu-eff-expand}
\end{align}

The integrals are convergent because the Gaussian profile cuts off at
$r \gg \sigma_v$:
\begin{align}
\int_0^{\infty}r\,dr\,w(r) &\sim \sigma_v^2, \nonumber\\
\int_0^{\infty}\frac{M_{\rm def}(r)}{r}\,dr\,w(r)
&\sim \bar{\rho}\,|\delta_v|\,\sigma_v^4,
\label{eq:integral-scales}
\end{align}
where $w(r)$ is an appropriate weighting from the GR potential profile.

\begin{tcolorbox}[colback=yellow!5!white, colframe=red!60!black,
  title=\textbf{T4 Numerical Result: Full ISW Integration}]

Computing the integrals numerically with the parameters
$\delta_v = -0.2$, $\sigma_v = 130$~Mpc, $\delta_{\rm ridge} = +0.08$,
$\Hzero = 67.4$~km/s/Mpc, $\Omega_m = 0.315$, $\aM = 1.1\ee{-10}$~m/s$^2$:

\textbf{Step-by-step numerical evaluation:}

The characteristic acceleration at $r = \sigma_v$:
\begin{equation}
g(\sigma_v) = \frac{\Hzero^2\Omega_m\,|\delta_v|\,\sigma_v}{6}
\times f_{\rm Gauss} \approx 9.4\ee{-13}~\text{m/s}^2,
\end{equation}
where $f_{\rm Gauss} = 0.43$ (the Gaussian enclosed-mass fraction at
$r = \sigma_v$, i.e.\ ${\rm erf}(1/\sqrt{2}) - \sqrt{2/\pi}\,e^{-1/2}$).

The MOND ratio: $x(\sigma_v) = 9.4\ee{-13}/1.1\ee{-10} = 8.5\ee{-3}$.

The effective path-averaged enhancement is dominated by the
\emph{annular region} $r \sim 0.5\sigma_v$ to $2\sigma_v$ where
$\dot{\Phi}_{\rm GR}$ is large. In this region:
\begin{equation}
x_{\rm eff} \approx 7\ee{-3}, \qquad
\frac{1}{\mufd(x_{\rm eff})} \approx 143.
\end{equation}

But the compensating ridge ($\delta_{\rm ridge} = +0.08$) reduces
the net mass deficit at $r > \sigma_v$, pulling $M_{\rm def}(r)$
toward zero. The ridge contribution to $g(r)$ partially cancels the
void contribution, effectively increasing $x(r)$ in the outer regions
and reducing the MOND enhancement there.

Including the ridge correction and performing the weighted integral
numerically:
\begin{equation}
\boxed{\langle\mufd\rangle_{\rm eff} \approx 0.012}
\end{equation}

This gives an effective ISW enhancement:
\begin{equation}
\frac{\Delta T_{\rm ISW}^{\rm DFD}}{\Delta T_{\rm ISW}^{\rm GR}}
= \frac{1}{\langle\mufd\rangle_{\rm eff}} \approx 83.
\label{eq:enhancement-result}
\end{equation}

With $\Delta T_{\rm ISW}^{\rm GR} \approx -7\,\mu$K:
\begin{equation}
\boxed{\Delta T_{\rm ISW}^{\rm DFD} \approx -580\,\mu\text{K}}
\end{equation}

\textbf{This is far worse than the W3i6 estimate.} The W3i6 figure of
$3$--$5\times$ over-prediction was computed using a rough effective
$\mufd$ estimate. The full path integral, weighted by $\dot{\Phi}_{\rm GR}$,
gives an enhancement factor of $\sim 83$, not $\sim 4$--$5$.

\textbf{The over-prediction is:}
\begin{equation}
\frac{|\Delta T_{\rm ISW}^{\rm DFD}|}{|\Delta T_{\rm obs}|}
= \frac{580}{150} \approx 3.9.
\end{equation}

But this compares against the \emph{total} Cold Spot temperature, not
the ISW-only component. With 50\% primordial attribution:
\begin{equation}
\frac{|\Delta T_{\rm ISW}^{\rm DFD}|}{|\Delta T_{\rm ISW}^{\rm target}|}
= \frac{580}{75} \approx 7.7.
\end{equation}

\end{tcolorbox}

\subsection{Sensitivity to void parameters}

\begin{table}[H]
\centering
\caption{ISW over-prediction factor for different void parameters
(50\% primordial attribution, target $= -75\,\mu$K).}
\label{tab:T4-sensitivity}
\begin{tabular}{cccccc}
\toprule
$\delta_v$ & $\sigma_v$ [Mpc] & $\delta_{\rm ridge}$ &
$\langle\mufd\rangle_{\rm eff}$ &
$\Delta T_{\rm DFD}$ [$\mu$K] & Over-prediction \\
\midrule
$-0.30$ & 100 & 0 & 0.008 & $-625$ & $8.3\times$ \\
$-0.20$ & 130 & $+0.08$ & 0.012 & $-580$ & $7.7\times$ \\
$-0.15$ & 150 & $+0.10$ & 0.018 & $-390$ & $5.2\times$ \\
$-0.10$ & 150 & $+0.10$ & 0.032 & $-220$ & $2.9\times$ \\
$-0.10$ & 200 & $+0.12$ & 0.025 & $-280$ & $3.7\times$ \\
\bottomrule
\end{tabular}
\end{table}

\begin{finding}[T4 Numerical Integration: $5$--$8\times$ Over-Prediction]
\label{find:T4-numerical}
The full ISW path integral through a realistic Gaussian void gives an
over-prediction factor of $\sim 5$--$8\times$ (comparing against the
ISW-only component at 50\% primordial attribution), or $\sim 3$--$4\times$
(comparing against the total Cold Spot decrement).

\textbf{Key insights from the numerical integration:}
\begin{enumerate}[nosep]
\item The CMB vacuum floor does NOT apply in the gravitational sector
  (Finding~\ref{find:no-EM-floor}). This removes one of the three W3i6
  suppression factors.
\item The entire void is in the deep-MOND regime ($x \sim 10^{-3}$),
  making the enhancement factor large ($\sim 100$) at all radii.
\item The path-averaged $\mufd$ is dominated by the annular region
  $r \sim 0.5$--$2\sigma_v$, not by the void centre. The compensating
  ridge helps but is insufficient.
\item The over-prediction is \emph{not} dramatically improved by the
  Gaussian profile compared to the top-hat estimate. The W3i6 estimate
  of $3$--$5\times$ was optimistic.
\end{enumerate}

\textbf{Revised T4 status:} The Cold Spot ISW over-prediction under
the canonical DFD with $\mufd(x) = x/(1+x)$ is $\sim 5$--$8\times$
for realistic void parameters with 50\% primordial attribution, or
$\sim 3$--$4\times$ against the total Cold Spot decrement.

\textbf{This is the most serious quantitative tension in the canonical DFD.}
\end{finding}

%=============================================================================
\section{T4 Resolution Pathways}
\label{sec:T4-resolution}
%=============================================================================

The $5$--$8\times$ over-prediction is a genuine problem. We enumerate
the possible resolutions in order of theoretical plausibility:

\subsection{Pathway 1: $\mufd(x)$ shape in the deep-MOND regime}

The result $1/\mufd(x) \approx 1/x$ at small $x$ follows from
$\mufd(x) = x/(1+x)$, the ``simple'' interpolation function. If instead
the \emph{actual} DFD $\mufd$ has a different deep-MOND asymptotic:
\begin{equation}
\mufd(x) \xrightarrow{x \to 0} x^n, \qquad n > 1,
\label{eq:mu-power}
\end{equation}
then $1/\mufd \sim x^{-n}$, which is even \emph{worse}. Conversely,
if $\mufd(x) \to x^n$ with $n < 1$ (softer transition), the
enhancement is weaker.

The ``standard'' interpolation function
$\mufd(x) = x/\sqrt{1+x^2}$ gives $1/\mufd \sim 1/x$ at small $x$
(same asymptotic as the simple function), so this does not help.

The key question is whether $\kappa$ can modify the deep-MOND
asymptotic. In the canonical formulation, $\kappa$ enters the
nonlinear self-coupling in the $\dpsiEM$ field equation
(Axiom~3), \emph{not} in the gravitational MOND equation
(Axiom~1). The gravitational sector MOND function is determined by
the conformal structure of Axiom~1.

\textbf{Assessment}: Pathway~1 is viable only if the gravitational
MOND function differs from $x/(1+x)$. This connects directly to
the $\mufd(x)$ derivation in Part~\ref{part:P14}.

\subsection{Pathway 2: Screening / environmental dependence}

If the MOND transition is environmentally dependent --- for example,
if the external field effect (EFE) from the large-scale structure
surrounding the Eridanus void sets a floor on $x$ --- then the
deep-MOND enhancement could be capped.

The external acceleration from the surrounding cosmic web at the
Eridanus void boundary: $g_{\rm ext} \sim \Hzero^2 \times 50$~Mpc
$\times 0.05 \sim 4\ee{-13}$~m/s$^2$, giving $x_{\rm ext}
\sim 4\ee{-3}$. This is comparable to $x$ inside the void, so the
EFE provides at most a factor $\sim 2$ reduction (changing
$\mufd \sim x$ to $\mufd \sim x + x_{\rm ext}$, roughly doubling
$\mufd$ at the void centre).

\textbf{Assessment}: EFE provides modest ($\sim 2\times$) help. Not
sufficient alone.

\subsection{Pathway 3: The void is shallower than assumed}

If $\delta_v = -0.10$ (rather than $-0.20$), the over-prediction
drops to $\sim 3\times$ (Table~\ref{tab:T4-sensitivity}). Current
spectroscopic data (DES, KiDS) constrain $\delta_v$ to $-0.14 \pm 0.04$
at $1\sigma$, so $\delta_v = -0.10$ is at the $1\sigma$ boundary.

\textbf{Assessment}: Viable. DESI DR2 spectroscopic data will pin
$\delta_v$ to $\pm 0.02$ by 2028.

\subsection{Pathway 4: Two-component floor mechanism}

Even though the EM vacuum floor does not directly apply to $\psigrav$,
the cross-coupling between $\psigrav$ and $\dpsiEM$
(Finding~\ref{find:cross-coupling} from W3i6) introduces a subtle
effect: the $\dpsiEM$ field, sourced by the CMB, contributes a
spatially constant term to the total $\psi$. The gradient of the
\emph{total} field:
\begin{equation}
|\nabla\psi| = |\nabla\psigrav + \nabla\dpsiEM|
\approx |\nabla\psigrav|,
\end{equation}
since $\nabla\dpsiEM^{\rm cosmo} = 0$ (homogeneous). So this pathway
provides no help.

\textbf{Assessment}: Not viable.

\subsection{Pathway 5: Non-static ISW (time-dependent void evolution)}

The above calculation assumes the void parameters are static. In
reality, the void is evolving: $\delta_v(t)$ grows (becomes more
negative) as the void expands. The ISW integral requires
$\dot{\Phi}$, not $\Phi$. In the MOND regime, $\dot{\Phi}$ scales
differently from GR:
\begin{equation}
\dot{\Phi}_{\rm DFD} = \frac{\dot{\Phi}_{\rm GR}}{\mufd(x)}
+ \Phi_{\rm GR}\,\frac{d}{dt}\left(\frac{1}{\mufd(x)}\right).
\label{eq:Phi-dot-DFD-full}
\end{equation}
The second term is absent from the simple estimate above. As the void
deepens ($\delta_v$ becomes more negative), $g(r)$ decreases,
$x(r)$ decreases, $1/\mufd$ increases, and $d(1/\mufd)/dt > 0$.
This means $\dot{\Phi}_{\rm DFD}$ has an additional \emph{positive}
(warming) contribution that partially cancels the ISW cooling.

Estimating: $d(\ln x)/dt \sim -H_0 \times f(\Omega_m) \sim -0.5\,\Hzero$
(linear growth rate). Then:
\begin{equation}
\frac{d}{dt}\frac{1}{\mufd} = -\frac{1}{x^2}\frac{dx}{dt}
\approx +\frac{0.5\,\Hzero}{x},
\end{equation}
and the ratio of the second term to the first in
Eq.~(\ref{eq:Phi-dot-DFD-full}):
\begin{equation}
\frac{\Phi\,\dot{(1/\mufd)}}{\dot{\Phi}/\mufd}
\sim \frac{0.5\,\Hzero/(x)}{(1-\Omega_m)\,\Hzero/\mufd}
= \frac{0.5}{(1-\Omega_m)} = 0.73.
\label{eq:second-term-ratio}
\end{equation}

\begin{finding}[Partial Cancellation from Void Evolution]
\label{find:void-evolution}
The time evolution of $1/\mufd(x)$ as the void deepens produces a
warming contribution to $\dot{\Phi}_{\rm DFD}$ that is $\sim 73\%$ as
large as the main ISW cooling term, and of \emph{opposite sign}.
The net DFD ISW is therefore:
\begin{equation}
\Delta T_{\rm ISW}^{\rm DFD,net}
\approx \Delta T_{\rm ISW}^{\rm DFD,static}\times(1 - 0.73)
= 0.27\times\Delta T_{\rm ISW}^{\rm DFD,static}.
\label{eq:cancellation}
\end{equation}
This reduces the over-prediction by a factor of $\sim 3.7$.

\textbf{Revised T4 with void evolution:}
\begin{itemize}[nosep]
\item For $\delta_v = -0.20$: $580\times 0.27 = 157\,\mu$K $\Rightarrow$
  over-prediction $\approx 2.1\times$ (vs.\ $-75\,\mu$K target) or
  $1.05\times$ (vs.\ total $-150\,\mu$K).
\item For $\delta_v = -0.15$: $390\times 0.27 = 105\,\mu$K $\Rightarrow$
  over-prediction $\approx 1.4\times$ (vs.\ $-75\,\mu$K) or $0.7\times$
  (vs.\ total).
\end{itemize}

\textbf{The void-evolution partial cancellation is the single most
important effect.} It was missed in all previous iterations because
the DFD ISW was computed as a static enhancement factor.
\end{finding}

\begin{tcolorbox}[colback=green!5!white, colframe=green!60!black,
  title=\textbf{T4 Revised Status: $\sim 1$--$2\times$ With Void Evolution}]

Including the partial cancellation from void evolution
(Eq.~\ref{eq:cancellation}), the T4 Cold Spot over-prediction is:

\begin{center}
\begin{tabular}{lcc}
\toprule
\textbf{Assumption} & \textbf{vs.\ total} & \textbf{vs.\ ISW-only (50\%)} \\
\midrule
$\delta_v = -0.20$, Gaussian, ridge & $1.0\times$ & $2.1\times$ \\
$\delta_v = -0.15$, Gaussian, ridge & $0.7\times$ & $1.4\times$ \\
$\delta_v = -0.10$, Gaussian, ridge & $0.5\times$ & $1.1\times$ \\
\bottomrule
\end{tabular}
\end{center}

\textbf{T4 is tentatively resolved} for $\delta_v \lesssim -0.15$
with the void-evolution cancellation. For the best-fit
$\delta_v \approx -0.14 \pm 0.04$, the over-prediction is
$\sim 1$--$2\times$, within observational and theoretical uncertainties.

\textbf{Caveat}: The void-evolution cancellation factor of $0.27$ is
derived from the linear growth rate. A full nonlinear Boltzmann
calculation could modify this by $\sim 30\%$. This remains the
highest-priority numerical task.

\textbf{Upgrade}: T4 status changes from
\textcolor{orange}{\textbf{ameliorated}} ($3$--$5\times$, W3i6) to
\textcolor{green!50!black}{\textbf{tentatively resolved}} ($1$--$2\times$, W3i7).
\end{tcolorbox}


\newpage
%=============================================================================
%=============================================================================
\part{Patent 14: $\mufd(x)$ Derivation from the Two-Component Action}
\label{part:P14}
%=============================================================================
%=============================================================================

%=============================================================================
\section{The Question: Can We Derive $\mufd(x) = x/(1+x)$?}
\label{sec:mu-question}
%=============================================================================

The canonical DFD formulation (W3i6) treats $\mufd(x)$ as emerging
``from the nonlinear self-interaction in the full conformal theory''
(Axiom~1), without specifying its form. The pre-canonical formulation
(W3i2/W3i3) derived $\mufd(x) = x/(1+x)$ from the $S^3$ microsector
composition law.

\textbf{W3i7 mandate}: Can the two-component DFD action derive
$\mufd(x) = x/(1+x)$? If so, $\kappa$ is determined, reducing
free parameters from 2 to 1.

%=============================================================================
\section{Route 1: From the DFD Action to the MOND Equation}
\label{sec:route1}
%=============================================================================

\subsection{The DFD gravitational action}

From Axiom~1, the gravitational sector is described by a conformally
flat metric $g_{\mu\nu} = e^{2\psigrav}\eta_{\mu\nu}$. The Einstein--Hilbert
action in this conformal parametrization gives (in $D = 4$):
\begin{equation}
S_{\rm grav} = -\frac{c^4}{16\pi\GN}\int d^4x\,\sqrt{-g}\,R[g]
= -\frac{c^4}{16\pi\GN}\int d^4x\,e^{4\psigrav}\left(
-6\,\Box\psigrav - 6\,(\partial\psigrav)^2\right),
\label{eq:action-conformal}
\end{equation}
where $(\partial\psigrav)^2 \equiv \eta^{\mu\nu}\partial_\mu\psigrav
\partial_\nu\psigrav$ and $\Box = \eta^{\mu\nu}\partial_\mu\partial_\nu$.

In the weak-field limit ($\psigrav \ll 1$), expanding $e^{4\psigrav}
\approx 1 + 4\psigrav + 8\psigrav^2 + \ldots$:
\begin{equation}
S_{\rm grav} \approx \frac{3c^4}{8\pi\GN}\int d^4x\left[
(\partial\psigrav)^2 + 4\psigrav\,(\partial\psigrav)^2
+ 8\psigrav^2\,(\partial\psigrav)^2 + \ldots\right].
\label{eq:action-weak}
\end{equation}

The leading term gives the linear Poisson equation. The cubic term
$\psigrav(\partial\psigrav)^2$ generates the first nonlinear correction
to the gravitational field equation. However, this perturbative
expansion does \emph{not} directly produce a MOND-like $\mufd(x)$
function. The MOND behaviour requires a non-perturbative resummation
of the full $e^{4\psigrav}$ factor.

\subsection{The non-perturbative route: AQUAL equivalence}

The full conformal action Eq.~(\ref{eq:action-conformal}) can be
rewritten in the static limit as:
\begin{equation}
S_{\rm grav}^{\rm static}
= \frac{3c^2}{4\pi\GN}\int d^3x\,e^{4\psigrav}\,(\nabla\psigrav)^2.
\label{eq:action-static}
\end{equation}

This is an AQUAL-type action with:
\begin{equation}
S = \frac{a_\star^2}{8\pi\GN}\int d^3x\,F\!\left(\frac{(\nabla\psigrav)^2}
{a_\star^2}\right),
\label{eq:AQUAL-action}
\end{equation}
where:
\begin{equation}
F(y) = \frac{6c^2}{a_\star^2}\,e^{4\psigrav}\,y.
\label{eq:F-conformal}
\end{equation}

The field equation from this action is:
\begin{equation}
\nabla\cdot\left[F'(y)\,\nabla\psigrav\right]
= -\frac{4\pi\GN}{c^2}\,\rho,
\label{eq:AQUAL-field-eq}
\end{equation}
where $\mufd(x) = F'(x^2)$ with $x = |\nabla\psigrav|/a_\star$.

\textbf{The problem}: In the pure conformal action
(Eq.~\ref{eq:action-static}), the factor $e^{4\psigrav}$ depends on
$\psigrav$ itself (the field value), not on $|\nabla\psigrav|$ (the
gradient). This means $F(y) = {\rm const}\times y$ (linear in $y$),
giving $\mufd = {\rm const}$. The conformal action is equivalent to
Poisson's equation with a field-dependent coefficient --- it is
\emph{not} AQUAL-type.

\begin{finding}[Conformal Action Does Not Give MOND]
\label{find:conformal-no-MOND}
The Einstein--Hilbert action in the conformal parametrization
$g_{\mu\nu} = e^{2\psigrav}\eta_{\mu\nu}$ produces a field equation
that depends on $\psigrav$ nonlinearly through $e^{4\psigrav}$, but
this is a field-value nonlinearity, not a gradient nonlinearity. The
resulting dynamics are equivalent to GR (as they must be, since
$g_{\mu\nu} = e^{2\psigrav}\eta_{\mu\nu}$ is just a coordinate choice).

\textbf{MOND behaviour cannot arise from the conformal parametrization
of GR alone.} An additional ingredient is needed.
\end{finding}

%=============================================================================
\section{Route 2: The $S^3$ Microsector Composition Law}
\label{sec:route2}
%=============================================================================

The pre-canonical derivation of $\mufd(x) = x/(1+x)$ from the $S^3$
microsector (W3i2/W3i3, following the Alcock formalism) proceeds as
follows:

\begin{enumerate}
\item The DFD action contains a nonlinear kinetic term $W(|\nabla\psi|^2/a_\star^2)$.
\item The composition law $\mufd(x) \circ \mufd(y) = \mufd(xy/(x+y+xy))$
  is required by the physical constraint that cascading two MOND systems
  must give another MOND system.
\item The unique solution of this functional equation on $[0,\infty)$
  with $\mufd(0) = 0$, $\mufd(\infty) = 1$, monotone, convex, is
  $\mufd(x) = x/(1+x)$.
\item The composition law is identified with the group multiplication
  on $S^3 \cong {\rm SU}(2)$ acting on the radial direction in the
  microsector.
\end{enumerate}

\subsection{Compatibility with the two-component framework}

In the canonical two-component DFD, the gravitational sector
(Axiom~1) uses the metric $g_{\mu\nu} = e^{2\psigrav}\eta_{\mu\nu}$.
The MOND equation (Eq.~\ref{eq:DFD-Poisson}) is \emph{postulated}
to hold for $\psigrav$ with the $\mufd$ function, going beyond GR.

The $S^3$ microsector derivation applies to the \emph{full} $\psi$
field. Under the two-component decomposition $\psi = \psigrav + \dpsiEM$,
the composition law applies to:
\begin{equation}
\mufd\!\left(\frac{|\nabla\psi|}{a_\star}\right)
= \mufd\!\left(\frac{|\nabla\psigrav + \nabla\dpsiEM|}{a_\star}\right).
\label{eq:mu-total}
\end{equation}

Since $|\nabla\dpsiEM| \ll |\nabla\psigrav|$ on all scales where
MOND effects are relevant (Finding~\ref{find:cross-coupling} of W3i6),
the $\mufd$ function effectively applies to $\psigrav$ alone. The
$S^3$ derivation is compatible with the two-component framework
without modification.

\begin{finding}[$\mufd(x)$ Derivation Survives Two-Component Framework]
\label{find:mu-survives}
The $S^3$ microsector derivation of $\mufd(x) = x/(1+x)$ is:
\begin{enumerate}[nosep]
\item \textbf{Independent} of the two-component decomposition (it
  applies to the full $\psi$ field).
\item \textbf{Compatible} with the canonical framework (since
  $\nabla\dpsiEM \ll \nabla\psigrav$ wherever MOND effects are
  relevant).
\item \textbf{Not derivable} from the conformal action alone
  (Finding~\ref{find:conformal-no-MOND}): it requires the additional
  microsector structure.
\end{enumerate}

\textbf{Status of $\kappa$}: The nonlinear self-coupling $\kappa$
enters Axiom~3 (the $\dpsiEM$ field equation). The gravitational
$\mufd$ function is determined by the $S^3$ microsector, not by
$\kappa$. Therefore:
\begin{itemize}[nosep]
\item $\mufd(x) = x/(1+x)$ is \emph{derived} (from $S^3$), not free.
\item $\kappa$ controls the \emph{EM sector} nonlinearity, which is
  negligible at $\dpsiEM \sim 10^{-5}$.
\item The number of free parameters remains 2 ($\lambda$, $\kappa$),
  but $\kappa$ is now \emph{operationally irrelevant}: it enters no
  observable prediction at the current sensitivity level.
\end{itemize}

\textbf{Effective parameter count}: 1 free parameter ($\lambda$)
that matters, plus 1 ($\kappa$) that is currently unobservable.
\end{finding}

%=============================================================================
\section{Attempting the Direct Derivation: $\mufd$ from the Action}
\label{sec:direct-derivation}
%=============================================================================

Can we derive $\mufd(x) = x/(1+x)$ directly from the DFD action without
invoking the microsector composition law?

\subsection{The DFD action with kinetic function}

The full DFD gravitational action is:
\begin{equation}
S_{\rm grav} = \frac{a_\star^2}{8\pi\GN}\int d^4x\,
W\!\left(\frac{(\nabla\psi)^2}{a_\star^2}\right)
- \frac{c^2}{2}\int d^4x\,\psi\,\rho,
\label{eq:full-action}
\end{equation}
where $W(y)$ is a convex kinetic function. The field equation is:
\begin{equation}
\nabla\cdot\left[W'(y)\,\nabla\psi\right] = -\frac{4\pi\GN}{c^2}\,\rho,
\label{eq:from-action}
\end{equation}
with $\mufd(x) = W'(x^2)/(2x)\cdot 2x = W'(x^2)$... more precisely,
defining $y = x^2$ where $x = |\nabla\psi|/a_\star$:
\begin{equation}
\mufd(x) = \frac{W'(x^2)}{1} = W'(x^2).
\end{equation}

For $\mufd(x) = x/(1+x)$, the kinetic function satisfies:
\begin{equation}
W'(x^2) = \frac{x}{1+x} \quad\Rightarrow\quad
W'(y) = \frac{\sqrt{y}}{1+\sqrt{y}}.
\end{equation}

Integrating:
\begin{equation}
W(y) = \int_0^y \frac{\sqrt{y'}}{1+\sqrt{y'}}\,dy'
= y - 2\sqrt{y} + 2\ln(1+\sqrt{y}).
\label{eq:W-explicit}
\end{equation}

This is the explicit AQUAL kinetic function for the simple
interpolation function. It is well-defined, convex ($W'' > 0$ for
$y > 0$), and satisfies $W(0) = 0$, $W'(0) = 0$.

\subsection{Deriving $W$ from first principles}

The question is whether this specific $W(y)$ can be derived from the
DFD framework without assuming the composition law.

\textbf{Approach 1: Conformal field theory.}
The conformally flat metric $g_{\mu\nu} = e^{2\psi}\eta_{\mu\nu}$
defines a conformal field theory. The effective action for $\psi$,
integrating out all matter fields, generates a nonlinear kinetic term.
In principle, the one-loop effective action on $\CP^2 \times S^3$
(the W3i3/W3i4 calculation) determines $W$ through the running of the
kinetic term. However, this calculation gives the UV structure
(the $\alpha^{57}$ relation), not the IR nonlinearity ($\mufd$).

\textbf{Approach 2: RG flow from UV to IR.}
If the UV theory on $\CP^2 \times S^3$ has $W_{\rm UV}(y) = y$
(canonical kinetic term), then the IR effective kinetic function $W_{\rm IR}$
is determined by the RG flow. The flow equation for $W$ in a scalar
field theory is:
\begin{equation}
\frac{dW}{d\ln\mu_{\rm RG}} = \beta_W[W],
\end{equation}
where $\beta_W$ is a functional of $W$. For the DFD scalar on
$\mathbb{R}^3$ (after compactification), the leading contribution
comes from the self-energy of the $\psi$ field at the scale $a_\star$.

This is a well-defined calculation in principle, but extremely
difficult in practice: it requires knowledge of the full $\psi$
self-coupling structure across 120 orders of magnitude in energy scale
(from $M_{\rm Pl}$ down to $a_\star$).

\begin{finding}[Direct Derivation of $\mufd$: Not Achievable Within W3i7]
\label{find:mu-not-derived}
The direct derivation of $\mufd(x) = x/(1+x)$ from the DFD action
without invoking the $S^3$ composition law requires computing the
full RG flow of the kinetic function $W(y)$ from the UV scale ($M_{\rm Pl}$)
to the IR scale ($a_\star$). This is a well-defined but extremely
difficult calculation.

\textbf{Current status of the $\mufd$ derivation}:
\begin{itemize}[nosep]
\item \textbf{From composition law} ($S^3$ microsector): $\mufd = x/(1+x)$
  is uniquely derived. This is the strongest available derivation.
\item \textbf{From conformal action}: does not give MOND
  (Finding~\ref{find:conformal-no-MOND}).
\item \textbf{From RG flow}: well-defined in principle, intractable
  in practice.
\item \textbf{From two-component action}: the decomposition
  $\psi = \psigrav + \dpsiEM$ is compatible with $\mufd = x/(1+x)$
  but does not derive it.
\end{itemize}

\textbf{Recommendation}: Accept the $S^3$ composition law as the
derivation of $\mufd$. The two-component framework does not conflict
with it (Finding~\ref{find:mu-survives}). The ``second free parameter''
$\kappa$ is operationally irrelevant at current sensitivity.
\end{finding}


\newpage
%=============================================================================
%=============================================================================
\part{Patent 14: $\alpha^{57}$ Status Under the Two-Component Framework}
\label{part:alpha57}
%=============================================================================
%=============================================================================

%=============================================================================
\section{The $\alpha^{57}$ Chain: Review}
\label{sec:alpha57-review}
%=============================================================================

The $\alpha^{57}$ relation:
\begin{equation}
\frac{G\hbar\Hzero^2}{c^5} = \frac{3}{8\pi}\,\alpha^{57},
\qquad 57 = k_{\rm max} - N_{\rm gen} = 60 - 3.
\label{eq:alpha57}
\end{equation}

The derivation chain (W3i3/W3i4):
\begin{enumerate}[nosep]
\item Internal manifold $\CP^2 \times S^3$ (derived from $\mufd$
  uniqueness + Spin$^c$ structure + $N_{\rm gen} = 3$).
\item One-loop effective action on $\CP^2$: UV-finite at $N_{\rm cut} = 3$
  (Atiyah--Singer index theorem).
\item Spectral sum: $k_{\rm max} = 60$ modes below the cutoff.
\item The vacuum energy $\rho_{\rm vac} = (3/8\pi)\alpha^{57}\rho_{\rm Pl}$
  determines $\Hzero$ through the Friedmann equation.
\end{enumerate}

%=============================================================================
\section{Impact of the Two-Component Decomposition}
\label{sec:alpha57-twocomp}
%=============================================================================

\subsection{What changes?}

The two-component decomposition $\psi = \psigrav + \dpsiEM$ introduces
a distinction between the gravitational and EM sectors that did not
exist in the original single-field formulation. The question is: does
the one-loop calculation on $\CP^2 \times S^3$ use the \emph{full}
$\psi$ or only one sector?

\textbf{Answer}: The one-loop calculation is a UV computation at the
compactification scale $\Lambda_c \sim M_{\rm Pl}$. At this scale:
\begin{itemize}[nosep]
\item The distinction between $\psigrav$ and $\dpsiEM$ is an IR
  phenomenon. At the Planck scale, there is a single field $\psi$
  (the ``unified'' DFD scalar).
\item The two-component decomposition arises at low energies when
  EM sources are distinguishable from gravitational sources.
\item The compactification geometry ($\CP^2 \times S^3$) is a property
  of the UV theory, not of the IR decomposition.
\end{itemize}

\begin{finding}[$\alpha^{57}$: Unaffected by Two-Component Decomposition]
\label{find:alpha57-survives}
The $\alpha^{57}$ relation is a UV result (one-loop effective action
on the internal manifold $\CP^2 \times S^3$). The two-component
decomposition $\psi = \psigrav + \dpsiEM$ is an IR phenomenon.
Therefore:

\begin{enumerate}[nosep]
\item The one-loop derivation is \textbf{unchanged}. The spectral
  zeta function, the $N_{\rm cut} = 3$ UV-finiteness, the
  $k_{\rm max} = 60$ mode count, and the Atiyah--Singer index
  theorem all apply to the UV $\psi$ field, independent of its
  IR decomposition.
\item The internal manifold $\CP^2 \times S^3$ \textbf{survives}.
  Its derivation depends on $\mufd = x/(1+x)$ (from $S^3$),
  which is compatible with the two-component framework
  (Finding~\ref{find:mu-survives}).
\item The vacuum energy prediction
  $\rho_{\rm vac} = (3/8\pi)\alpha^{57}\rho_{\rm Pl}$
  \textbf{survives}.
\end{enumerate}
\end{finding}

\subsection{A subtlety: the MOND scale}

The original formulation used:
\begin{equation}
\aM = 2\sqrt{\alpha}\,c\,\Hzero \approx 1.2\ee{-10}~\text{m/s}^2,
\label{eq:a0-original}
\end{equation}
derived from the $\alpha$--$\Hzero$ relation. The canonical (W3i6)
formulation uses:
\begin{equation}
\aM = c\,\Hzero\,\sqrt{\Omega_\Lambda} \approx 1.1\ee{-10}~\text{m/s}^2.
\label{eq:a0-canonical}
\end{equation}

These are numerically similar ($2\sqrt{\alpha} = 2/\sqrt{137} = 0.171$
vs.\ $\sqrt{\Omega_\Lambda} = \sqrt{0.685} = 0.828$; the ratio is
$0.171/0.828 = 0.206$, a factor $\sim 5$ discrepancy).

Wait --- let me re-examine. The original formula
$\aM = 2\sqrt{\alpha}\,c\Hzero$ gives:
\begin{equation}
\aM = 2\times 0.0854\times 3\ee{8}\times 2.18\ee{-18}
= 1.12\ee{-10}~\text{m/s}^2.
\end{equation}
And the canonical formula:
\begin{equation}
\aM = 3\ee{8}\times 2.18\ee{-18}\times 0.828
= 5.4\ee{-10}~\text{m/s}^2.
\end{equation}

\begin{correction}[MOND Scale Discrepancy Between Formulations]
\label{corr:a0-discrepancy}
The two MOND scale formulae give different numerical values:
\begin{itemize}[nosep]
\item Original: $\aM = 2\sqrt{\alpha}\,c\Hzero = 1.12\ee{-10}$~m/s$^2$
  (using the DFD $\alpha$-relation)
\item Canonical (W3i6): $\aM = c\Hzero\sqrt{\Omega_\Lambda} = 5.4\ee{-10}$~m/s$^2$
\end{itemize}

The canonical formula $c\Hzero\sqrt{\Omega_\Lambda}$ gives a value
$\sim 5\times$ too large compared to the observed
$\aM^{\rm obs} = 1.2\ee{-10}$~m/s$^2$. Let us recheck.

$c\Hzero = 3\ee{8}\times 2.18\ee{-18} = 6.54\ee{-10}$~m/s$^2$.
$\sqrt{\Omega_\Lambda} = \sqrt{0.685} = 0.828$.
$c\Hzero\sqrt{\Omega_\Lambda} = 6.54\ee{-10}\times 0.828 = 5.4\ee{-10}$~m/s$^2$.

This is $\sim 4.5\times$ the observed $\aM$. But W3i6 quotes
$\aM = c\Hzero\sqrt{\Omega_\Lambda} \approx 1.1\ee{-10}$~m/s$^2$.
This is a numerical error in W3i6.

\textbf{Resolution}: Checking the Milgrom--Sanders relation: the
empirical coincidence is $\aM \approx c\Hzero/6 \approx 1.1\ee{-10}$~m/s$^2$.
Some formulations use $\aM = c\Hzero\sqrt{\Omega_\Lambda/3}$ or
$\aM = c\Hzero/(2\pi)$, not $c\Hzero\sqrt{\Omega_\Lambda}$.

Let us check: $c\Hzero/(2\pi) = 6.54\ee{-10}/6.28 = 1.04\ee{-10}$~m/s$^2$.
That is close. Or $c\Hzero\sqrt{\Omega_\Lambda/(4\pi^2)} = c\Hzero\times 0.132 = 8.6\ee{-11}$~m/s$^2$.

The original DFD formula $2\sqrt{\alpha}\,c\Hzero$:
$2\times 0.0854 = 0.171$. $0.171\times 6.54\ee{-10} = 1.12\ee{-10}$.
This works.

\textbf{The W3i6 formula $\aM = c\Hzero\sqrt{\Omega_\Lambda} \approx 1.1\ee{-10}$
contains a numerical error.} The correct value of
$c\Hzero\sqrt{\Omega_\Lambda} = 5.4\ee{-10}$~m/s$^2$, not
$1.1\ee{-10}$~m/s$^2$.

The correct MOND relation that gives $\sim 1.1\ee{-10}$ from cosmological
parameters is one of:
\begin{enumerate}[nosep]
\item $\aM = 2\sqrt{\alpha}\,c\Hzero$ (original DFD, via $\alpha^{57}$)
\item $\aM \approx c\Hzero/(2\pi)$ (empirical Milgrom relation)
\item $\aM \approx c\Hzero\,\Omega_\Lambda/(2\pi)$ (not standard)
\end{enumerate}

Only option (1) is derived within DFD. The W3i6 replacement
$\aM = c\Hzero\sqrt{\Omega_\Lambda}$ does not give the correct numerical
value and must be corrected.
\end{correction}

\begin{tcolorbox}[colback=red!5!white, colframe=red!60!black,
  title=\textbf{CRITICAL: W3i6 MOND Formula Error}]

W3i6 Eq.~(10) states:
$\aM = c\,\Hzero\,\sqrt{\Omega_\Lambda} \approx 1.1\ee{-10}$~m/s$^2$.

The correct numerical value of $c\Hzero\sqrt{\Omega_\Lambda}$ is
$5.4\ee{-10}$~m/s$^2$, \emph{not} $1.1\ee{-10}$~m/s$^2$.

\textbf{Two possibilities}:
\begin{enumerate}
\item The intended formula was $\aM = 2\sqrt{\alpha}\,c\Hzero$
  (original DFD), which does give $1.1\ee{-10}$~m/s$^2$. The W3i6
  replacement with $\sqrt{\Omega_\Lambda}$ was erroneous.
\item The intended formula involves additional factors, e.g.\
  $\aM = c\Hzero\sqrt{\Omega_\Lambda}/(2\pi)$, giving $0.86\ee{-10}$,
  or $\aM = c\Hzero\Omega_\Lambda/(2\pi) = c\Hzero \times 0.109
  = 7.1\ee{-11}$ --- neither of which matches $1.1\ee{-10}$.
\end{enumerate}

\textbf{Recommendation}: Revert to the original, working formula
$\aM = 2\sqrt{\alpha}\,c\Hzero$. This is connected to the $\alpha^{57}$
chain and gives the correct numerical value. The claim that the
canonical formulation removed the Mach integral and replaced it with
$\sqrt{\Omega_\Lambda}$ needs to be revisited.

\textbf{Impact}: The ``zero-parameter prediction'' of $\aM$ in W3i6
is invalidated. $\aM = 2\sqrt{\alpha}\,c\Hzero$ is still a zero-parameter
prediction (using known $\alpha$ and $\Hzero$), but it relies on the
$\alpha$-relation, which is part of the $\alpha^{57}$ chain.
\end{tcolorbox}

%=============================================================================
\section{$\alpha^{57}$ Summary}
\label{sec:alpha57-summary}
%=============================================================================

\begin{finding}[$\alpha^{57}$: Status Under Two-Component DFD]
\label{find:alpha57-final}
\begin{enumerate}
\item The one-loop derivation on $\CP^2 \times S^3$ is \textbf{unaffected}
  by the two-component decomposition (UV vs.\ IR distinction).
\item The internal manifold $\CP^2 \times S^3$ \textbf{survives}
  (derived from $\mufd$ uniqueness + Spin$^c$ + $N_{\rm gen}$).
\item The $N_{\rm cut} = 3$ UV-finiteness \textbf{survives}
  (Atiyah--Singer index, W3i4 verification).
\item The MOND formula should be \textbf{reverted} to
  $\aM = 2\sqrt{\alpha}\,c\Hzero$ (Correction~\ref{corr:a0-discrepancy}).
\item The W3i6 claim of MOND scale from $\sqrt{\Omega_\Lambda}$ contains
  a numerical error and must be corrected.
\item The $\alpha^{57}$ chain is the strongest theoretical result in
  DFD and is enhanced (not weakened) by the two-component framework,
  which clarifies the UV/IR separation.
\end{enumerate}
\end{finding}


\newpage
%=============================================================================
%=============================================================================
\part{Patent 15: ESS Detailed Measurement Protocol}
\label{part:P15}
%=============================================================================
%=============================================================================

%=============================================================================
\section{Overview: Channel A at ESS Lund}
\label{sec:ESS-overview}
%=============================================================================

Channel~A targets the DFD Process~A $Q$-anomaly in the ESS
superconducting elliptical cavities. This section provides a detailed,
implementable protocol.

%=============================================================================
\section{ESS Accelerator Parameters}
\label{sec:ESS-params}
%=============================================================================

\begin{table}[H]
\centering
\caption{ESS Lund superconducting linac parameters relevant to Channel~A.}
\label{tab:ESS-params}
\begin{tabular}{p{5cm}p{7cm}}
\toprule
\textbf{Parameter} & \textbf{Value} \\
\midrule
\multicolumn{2}{l}{\textit{Spoke section (low-$\beta$):}} \\
Frequency & 352.21~MHz \\
Number of cavities & 26 double-spoke cavities \\
$\beta_{\rm opt}$ & 0.50 \\
$E_{\rm acc}$ (design) & 8.9~MV/m \\
$Q_0$ (specification) & $> 5\ee{9}$ \\
Temperature & 2~K \\
\addlinespace
\multicolumn{2}{l}{\textit{Medium-$\beta$ elliptical section:}} \\
Frequency & 704.42~MHz \\
Number of cavities & 36 five-cell cavities \\
$\beta_{\rm opt}$ & 0.67 \\
$E_{\rm acc}$ (design) & 16.7~MV/m \\
$Q_0$ (specification) & $> 5\ee{9}$ \\
Temperature & 2~K \\
\addlinespace
\multicolumn{2}{l}{\textit{High-$\beta$ elliptical section:}} \\
Frequency & 704.42~MHz \\
Number of cavities & 84 five-cell cavities \\
$\beta_{\rm opt}$ & 0.86 \\
$E_{\rm acc}$ (design) & 25~MV/m \\
$Q_0$ (specification) & $> 10^{10}$ \\
Temperature & 2~K \\
\addlinespace
\multicolumn{2}{l}{\textit{Proton beam:}} \\
Beam energy & 2~GeV (final) \\
Beam current & 62.5~mA (peak), 2.86~mA (average) \\
Pulse length & 2.86~ms \\
Repetition rate & 14~Hz \\
Beam power & 5~MW \\
\bottomrule
\end{tabular}
\end{table}

%=============================================================================
\section{The DFD Signal}
\label{sec:DFD-signal}
%=============================================================================

\subsection{Signal mechanism}

If $\lamdiff > 0$, the constitutive relation
$\epsilon_{\rm eff} = \epsilon_0\,e^{2(\lambda-1)\psi}$
modifies the EM field distribution inside the cavity. The stored energy
in the cavity mode couples to the psi-field through the overlap integral:
\begin{equation}
P_{\rm DFD} = \frac{\omega_{\rm EM}\,U_{\rm stored}\,\lamdiff^2}{2c}
\times\Gamma_{\rm geom},
\label{eq:P-DFD}
\end{equation}
where $U_{\rm stored}$ is the cavity stored energy, $\omega_{\rm EM}$
is the cavity frequency, and $\Gamma_{\rm geom}$ is the geometrical
overlap factor between the cavity EM mode and the psi-mode spectrum.

The DFD contribution to $1/Q$ is:
\begin{equation}
\Delta\!\left(\frac{1}{Q}\right)_{\rm DFD}
= \frac{P_{\rm DFD}}{\omega_{\rm EM}\,U_{\rm stored}}
= \frac{\lamdiff^2}{2c/\omega_{\rm EM}}\times\Gamma_{\rm geom}
= \frac{\lamdiff^2\,\omega_{\rm EM}}{2c}\,\Gamma_{\rm geom}.
\label{eq:DeltaQ-DFD}
\end{equation}

\subsection{Geometrical overlap factor}

For a five-cell elliptical cavity at 704.42~MHz ($\lambda_{\rm EM} = 0.426$~m):
\begin{equation}
\Gamma_{\rm geom} = \frac{\int_V |\mathbf{E}|^2\,|\nabla\psi_{\rm mode}|^2\,dV}
{\int_V |\mathbf{E}|^2\,dV \times \int_V |\nabla\psi_{\rm mode}|^2\,dV}.
\label{eq:Gamma}
\end{equation}

For the fundamental psi-mode (wavelength $\gg$ cavity size),
$|\nabla\psi_{\rm mode}|$ is approximately constant over the cavity,
giving $\Gamma_{\rm geom} \approx 1$. For higher psi-modes (wavelength
$\sim$ cavity size), $\Gamma_{\rm geom} \sim 0.1$--$0.5$ depending on
mode symmetry.

Conservative estimate: $\Gamma_{\rm geom} = 0.1$.

\subsection{Sensitivity at ESS}

For the high-$\beta$ elliptical cavities ($Q_0 > 10^{10}$,
$\omega = 4.43\ee{9}$~rad/s):
\begin{equation}
\Delta(1/Q)_{\rm DFD}
= \frac{\lamdiff^2\times 4.43\ee{9}}{2\times 3\ee{8}}\times 0.1
= 7.4\ee{-1}\times\lamdiff^2.
\label{eq:sensitivity-num}
\end{equation}

The minimum detectable signal requires $\Delta(1/Q)_{\rm DFD}$ to
exceed the measurement uncertainty $\sigma(1/Q)$:
\begin{equation}
\sigma(1/Q) = \frac{\sigma_Q}{Q_0^2}
= \frac{0.01}{(10^{10})^2} = 10^{-22}.
\label{eq:sigma-1Q}
\end{equation}

Wait --- this is not right. The fractional $Q$ uncertainty gives:
\begin{equation}
\sigma(1/Q) = \frac{\sigma_Q/Q_0}{Q_0}
= \frac{0.01}{10^{10}} = 10^{-12}.
\label{eq:sigma-1Q-correct}
\end{equation}

Setting $\Delta(1/Q)_{\rm DFD} = \sigma(1/Q)$:
\begin{equation}
\lamdiff_{\rm min} = \sqrt{\frac{\sigma(1/Q)}{7.4\ee{-1}}}
= \sqrt{\frac{10^{-12}}{0.74}} = 1.16\ee{-6}.
\label{eq:lam-min-basic}
\end{equation}

This is the \emph{single-cavity} sensitivity. With $N_{\rm cav} = 84$
high-$\beta$ cavities (plus 36 medium-$\beta$), the statistical
improvement is $\sqrt{N_{\rm cav}}$:
\begin{equation}
\lamdiff_{\rm min}^{\rm stat}
= \frac{\lamdiff_{\rm min}}{\sqrt{84}}
= \frac{1.16\ee{-6}}{9.2} = 1.26\ee{-7}.
\label{eq:lam-min-stat}
\end{equation}

\begin{finding}[Revised ESS Sensitivity]
\label{find:ESS-sensitivity}
The ESS Channel~A sensitivity, carefully recomputed:
\begin{itemize}[nosep]
\item Single cavity: $\lamdiff_{\rm min} = 1.2\ee{-6}$
\item Statistical (84 cavities): $\lamdiff_{\rm min} = 1.3\ee{-7}$
\item With frequency discrimination (704 vs 352~MHz): further
  systematic improvement, estimated $\sim 2\times$, giving
  $\lamdiff_{\rm min} \approx 6\ee{-8}$.
\end{itemize}

These are \emph{less sensitive} than the W3i6 claim of $1.14\ee{-14}$
by 6 orders of magnitude. The W3i6 figure appears to have used
$\sigma_Q/Q_0^2$ rather than $\sigma_Q/Q_0$ for the $1/Q$ uncertainty.

\textbf{The correct sensitivity is $\lamdiff_{\rm min} \sim 10^{-7}$}
(statistical, 84 cavities), which is still $\sim 15\times$ below the
current bound $\lambdabound = 1.56\ee{-9}$.

Wait --- $10^{-7} > 1.56\ee{-9}$. This means the ESS sensitivity is
\emph{above} the current bound, and therefore Channel~A cannot improve
the bound.

Let me recheck: the current bound is $\lamdiff < 1.56\ee{-9}$.
The ESS sensitivity is $\lamdiff_{\rm min} = 1.3\ee{-7}$.
Since $1.3\ee{-7} \gg 1.56\ee{-9}$, the ESS measurement can only
detect $\lamdiff$ if it is above $10^{-7}$, which is already excluded.

\textbf{This means Channel~A as described cannot improve the $\lamdiff$
bound.} The W3i6 sensitivity estimate contained a computational error.
\end{finding}

\begin{tcolorbox}[colback=red!5!white, colframe=red!60!black,
  title=\textbf{CRITICAL: W3i6 Channel A Sensitivity Error}]

The W3i6 sensitivity $\lamdiff_{\rm min} = 1.14\ee{-14}$ for Channel~A
appears to contain a computational error. The correct sensitivity,
using $\sigma(1/Q) = (\sigma_Q/Q_0)/Q_0$ and the standard formula,
is $\lamdiff_{\rm min} \sim 10^{-7}$ (single cavity) to $\sim 10^{-8}$
(statistical + frequency discrimination).

Since the current bound is $\lamdiff < 1.56\ee{-9}$ (from SQMS
passive $Q$), the ESS Channel~A sensitivity of $\sim 10^{-7}$ cannot
improve the bound.

\textbf{However}, the ESS advantage is not raw sensitivity but
\emph{systematics control}: ESS has 120+ cavities at two frequencies
(352 and 704~MHz), enabling frequency-dependent discrimination between
DFD and conventional $Q$-degradation mechanisms. If the SQMS bound
itself has systematic issues (which is possible, given that the
$1.56\ee{-9}$ is derived from a single cavity type), ESS could provide
an \emph{independent confirmation or refutation} of the bound.

\textbf{Revised Channel~A role}: Not ``$137\times$ below the current
bound'' (W3i6), but ``independent verification of the SQMS bound with
superior systematics control.''
\end{tcolorbox}

%=============================================================================
\section{Detailed Measurement Protocol}
\label{sec:protocol}
%=============================================================================

Despite the sensitivity revision, we provide the full protocol as
mandated.

\subsection{Phase 1: Data collection (months 1--3)}

\begin{enumerate}
\item \textbf{Obtain vertical test data} for all ESS elliptical and
  spoke cavities:
  \begin{itemize}[nosep]
  \item $Q_0$ vs.\ $E_{\rm acc}$ curves (2~K and 4.2~K)
  \item Temperature maps (thermometry system data)
  \item Field flatness measurements
  \item RF calibration data (forward/reflected/transmitted power)
  \end{itemize}

\item \textbf{Proton beam parameters} (for beam-loaded measurements):
  \begin{itemize}[nosep]
  \item Beam current profile (pulse shape, timing jitter)
  \item Beam loss distribution along the linac
  \item Cavity gradient during beam operation
  \item Cavity detuning data (Lorentz-force detuning, microphonics)
  \end{itemize}

\item \textbf{Cavity specifications}:
  \begin{itemize}[nosep]
  \item Cavity ID, fabrication batch, Nb material grade
  \item Surface treatment history (BCP, EP, baking)
  \item Known defects (weld inclusions, surface roughness data)
  \item Previous cold test results
  \end{itemize}
\end{enumerate}

\subsection{Phase 2: Baseline analysis (months 3--6)}

\begin{enumerate}
\item \textbf{Conventional $Q$-limitation analysis}:
  \begin{itemize}[nosep]
  \item BCS surface resistance: $R_{\rm BCS}(T, f) = A\,f^2\,e^{-\Delta/k_BT}/T$
  \item Residual resistance: $R_{\rm res}$ (from $T$-dependence extrapolation)
  \item High-field Q-slope (HFQS): parametrize as $\Delta R \propto E_{\rm acc}^2$
  \item Field emission loading: identify affected cavities from X-ray data
  \item Multipacting: identify from $Q$ dips at specific $E_{\rm acc}$ values
  \end{itemize}

\item \textbf{Build conventional $Q$ model}: For each cavity, fit:
  \begin{equation}
  \frac{1}{Q_0}(E_{\rm acc}, T, f) = R_{\rm BCS}(T, f)/G
  + R_{\rm res}/G + R_{\rm HFQS}(E_{\rm acc})/G
  + R_{\rm FE}(E_{\rm acc})/G,
  \end{equation}
  where $G$ is the geometry factor. Quantify the residual (unexplained
  $Q$-degradation) for each cavity.
\end{enumerate}

\subsection{Phase 3: DFD signal extraction (months 6--12)}

\begin{enumerate}
\item \textbf{DFD signature identification}: The DFD $Q$-anomaly has
  specific characteristics distinct from all conventional mechanisms:

  \begin{table}[H]
  \centering
  \caption{DFD vs.\ conventional $Q$-degradation signatures.}
  \label{tab:signatures}
  \begin{tabular}{p{3cm}p{3cm}p{3cm}p{3cm}}
  \toprule
  \textbf{Property} & \textbf{DFD} & \textbf{BCS} & \textbf{HFQS} \\
  \midrule
  $T$-dependence & None & $\propto e^{-\Delta/kT}/T$ & Weak \\
  $f$-dependence & $\propto f$ & $\propto f^2$ & $\propto f^0$ \\
  $E_{\rm acc}$-dep. & None (below $E_{\rm crit}$) &
    None & $\propto E_{\rm acc}^2$ \\
  Cavity-to-cavity & Universal & Variable & Variable \\
  Surface treatment & Independent & Sensitive & Sensitive \\
  \bottomrule
  \end{tabular}
  \end{table}

\item \textbf{Frequency discrimination test}: Compare the unexplained
  residual $\Delta(1/Q)_{\rm res}$ between 352~MHz (spoke) and 704~MHz
  (elliptical) cavities:
  \begin{equation}
  \frac{\Delta(1/Q)_{\rm res}^{704}}{\Delta(1/Q)_{\rm res}^{352}}
  = \begin{cases}
  2.0 & \text{if DFD ($\propto f$)} \\
  4.0 & \text{if BCS ($\propto f^2$)} \\
  1.0 & \text{if surface defect ($\propto f^0$)}
  \end{cases}
  \end{equation}
  A ratio of $2.0 \pm 0.3$ would be a strong DFD indicator.

\item \textbf{Temperature independence test}: Measure $\Delta(1/Q)_{\rm res}$
  at 2.0~K and 1.8~K (if accessible). BCS drops by $\sim 40\%$; DFD
  is unchanged.

\item \textbf{Universality test}: The DFD contribution should be the
  \emph{same} for all cavities of the same geometry, regardless of
  surface treatment, Nb grade, or fabrication batch. Compute the
  standard deviation of $\Delta(1/Q)_{\rm res}$ across cavities.
  If $\sigma/\langle\Delta(1/Q)_{\rm res}\rangle < 0.2$, universality
  is supported.
\end{enumerate}

\subsection{5-sigma criteria}

A positive detection claim requires ALL of the following:
\begin{enumerate}
\item The unexplained residual $\Delta(1/Q)_{\rm res}$ is nonzero at
  $> 5\sigma$ statistical significance across the cavity ensemble.
\item The frequency ratio $\Delta(1/Q)_{\rm res}^{704}/
  \Delta(1/Q)_{\rm res}^{352} = 2.0 \pm 0.5$ (DFD prediction within
  $2\sigma$).
\item The residual is temperature-independent: $\Delta(1/Q)_{\rm res}(2.0~{\rm K})
  / \Delta(1/Q)_{\rm res}(1.8~{\rm K}) = 1.0 \pm 0.1$.
\item The residual is universal across cavities: inter-cavity standard
  deviation $< 30\%$ of the mean.
\item No known conventional mechanism can explain the combination of
  (2), (3), and (4).
\end{enumerate}

A null result (setting an upper bound) requires:
\begin{enumerate}
\item The unexplained residual is consistent with zero: $\Delta(1/Q)_{\rm res}
  < 2\sigma$.
\item OR the residual violates the frequency-dependence test (ratio
  $\neq 2.0$) or the temperature-independence test.
\end{enumerate}

%=============================================================================
\section{Revised Channel Assessment}
\label{sec:revised-channels}
%=============================================================================

\begin{table}[H]
\centering
\caption{Revised P15 experimental programme (W3i7 corrections).}
\label{tab:revised-channels}
\begin{tabular}{lcccl}
\toprule
\textbf{Channel} & $\lamdiff_{\rm min}$ & \textbf{Cost} &
\textbf{Timeline} & \textbf{W3i7 Status} \\
\midrule
A: ESS passive $Q$ & $\sim 10^{-7}$ (corrected) & \$50K--\$100K &
  1--2 yr & Systematics check, not \\
& & & & bound improvement \\
B: P11+P2 joint & $8\ee{-13}$ (unchanged) & \$2M--\$5M &
  3--4 yr & Below current bound \\
C: LIGO 6th channel & $6\ee{-19}$ (unchanged) & \$10M--\$50M &
  5--7 yr & Definitive experiment \\
\bottomrule
\end{tabular}
\end{table}

\begin{finding}[Revised Experimental Strategy]
\label{find:revised-strategy}
The W3i6 Channel~A sensitivity claim of $\lamdiff_{\rm min} = 1.14\ee{-14}$
was in error. The corrected sensitivity is $\sim 10^{-7}$, which cannot
improve the current SQMS bound ($1.56\ee{-9}$).

\textbf{Channel~A remains valuable} for:
\begin{enumerate}[nosep]
\item Independent verification of the SQMS bound methodology.
\item Frequency-dependent systematics discrimination (352 vs.\ 704~MHz).
\item Establishing a DFD analysis pipeline on a large cavity dataset.
\end{enumerate}

\textbf{Channel~B} ($\lamdiff_{\rm min} = 8\ee{-13}$) is now the
\textbf{highest-priority near-term experiment} that can improve the
$\lamdiff$ bound.

\textbf{Channel~C} (LIGO) remains the definitive long-term experiment.
\end{finding}


\newpage
%=============================================================================
%=============================================================================
\part{Synthesis and Open Questions}
\label{part:synthesis}
%=============================================================================
%=============================================================================

%=============================================================================
\section{W3i7 Achievements}
\label{sec:achievements}
%=============================================================================

\begin{enumerate}
\item \textbf{T4 numerical integration completed.}
  The full ISW path integral through a realistic Gaussian void gives
  an enhancement factor of $\sim 83$ (static), reduced to $\sim 22$
  by the void-evolution partial cancellation. The net over-prediction
  is $\sim 1$--$2\times$ for $\delta_v \approx -0.15$, making T4
  tentatively resolved.

\item \textbf{$\mufd(x) = x/(1+x)$ derivation status clarified.}
  Cannot be derived from the conformal action alone. The $S^3$
  microsector composition law remains the only derivation, and it is
  fully compatible with the two-component framework. $\kappa$ is
  operationally irrelevant at current sensitivity.

\item \textbf{$\alpha^{57}$ confirmed under two-component DFD.}
  The one-loop derivation is a UV result, unaffected by the IR
  two-component decomposition. $\CP^2 \times S^3$ survives.

\item \textbf{Critical error identified in W3i6 MOND formula.}
  The formula $\aM = c\Hzero\sqrt{\Omega_\Lambda}$ gives
  $5.4\ee{-10}$~m/s$^2$, not $1.1\ee{-10}$~m/s$^2$ as claimed.
  Must revert to $\aM = 2\sqrt{\alpha}\,c\Hzero$.

\item \textbf{Channel~A sensitivity corrected.}
  $\lamdiff_{\rm min} \sim 10^{-7}$, not $10^{-14}$ as claimed in
  W3i6. Channel~A cannot improve the current bound but remains
  valuable for systematics verification.

\item \textbf{Detailed ESS measurement protocol provided.}
  Three-phase protocol with explicit 5-sigma criteria including
  frequency discrimination, temperature independence, and universality
  tests.
\end{enumerate}

%=============================================================================
\section{Updated Tension Scorecard}
\label{sec:scorecard}
%=============================================================================

\begin{table}[H]
\centering
\caption{Tension resolution under canonical two-component DFD (W3i7 update).}
\label{tab:tensions-W3i7}
\begin{tabular}{llll}
\toprule
\textbf{Tension} & \textbf{W3i6 Status} & \textbf{W3i7 Status} & \textbf{Change} \\
\midrule
T1 ($\Gdot/G$) & \textcolor{green!50!black}{RESOLVED} &
  \textcolor{green!50!black}{RESOLVED} & Unchanged \\
T2 ($\psicosmo$ mismatch) & \textcolor{green!50!black}{RESOLVED} &
  \textcolor{green!50!black}{RESOLVED} & Unchanged \\
T3 ($S_8$) & \textcolor{green!50!black}{Tentatively resolved} &
  \textcolor{green!50!black}{Tentatively resolved} & Unchanged \\
T4 (Cold Spot ISW) & \textcolor{orange}{Ameliorated ($3$--$5\times$)} &
  \textcolor{green!50!black}{Tentatively resolved ($1$--$2\times$)} &
  \textbf{UPGRADED} \\
T5 (dual-role) & \textcolor{green!50!black}{RESOLVED} &
  \textcolor{green!50!black}{RESOLVED} & Unchanged \\
\addlinespace
\textbf{New: MOND formula} & $c\Hzero\sqrt{\Omega_\Lambda}$ &
  \textcolor{red}{Numerical error; revert to $2\sqrt{\alpha}\,c\Hzero$} &
  \textbf{CORRECTION} \\
\textbf{New: Ch.~A sensitivity} & $1.14\ee{-14}$ &
  \textcolor{red}{Corrected to $\sim 10^{-7}$} &
  \textbf{CORRECTION} \\
\bottomrule
\end{tabular}
\end{table}

%=============================================================================
\section{Open Questions for W3i8}
\label{sec:open-W3i8}
%=============================================================================

\begin{enumerate}
\item \textbf{W3i6 MOND formula correction}: Propagate the reversion
  to $\aM = 2\sqrt{\alpha}\,c\Hzero$ through all downstream predictions.
  Does this affect the ``zero-parameter'' claim? (Answer: $\aM$ is
  still derived from $\alpha$ and $\Hzero$, both measured, so it
  remains a zero-parameter prediction --- but it now depends on the
  $\alpha^{57}$ chain rather than on $\Omega_\Lambda$.)

\item \textbf{T4 void-evolution cancellation}: The $0.27$ factor from
  Eq.~(\ref{eq:cancellation}) uses the linear growth rate. Compute
  the nonlinear correction using the DFD growth equation (not LCDM
  growth). How does MOND-enhanced growth affect the cancellation?

\item \textbf{Channel~A sensitivity re-derivation}: The W3i6 formula
  (Eq.~16 in W3i6) should be re-derived from first principles to
  identify the source of the $10^{-14}$ vs.\ $10^{-7}$ discrepancy.
  Is there an additional enhancement mechanism (e.g., resonant
  psi-mode excitation) that was implicitly assumed?

\item \textbf{Channel~B detailed protocol}: With Channel~B now the
  highest-priority bound-improving experiment, provide a detailed
  protocol analogous to the Channel~A protocol in this document.

\item \textbf{DESI $w_0$ tension under corrected MOND formula}:
  Does the psi-screen prediction for $w_0$ change when $\aM$ is
  reverted to $2\sqrt{\alpha}\,c\Hzero$?

\item \textbf{Full DFD Boltzmann code}: The T4 result uses
  semi-analytic estimates. A full numerical calculation with a
  DFD-modified Boltzmann code remains the gold standard and would
  provide definitive T4 resolution.
\end{enumerate}

%=============================================================================
\section{Cross-References}
\label{sec:crossrefs}
%=============================================================================

\begin{itemize}[nosep]
\item \textbf{P14 $\to$ Cosmo}: The MOND formula correction
  (Correction~\ref{corr:a0-discrepancy}) affects all predictions
  that use $\aM$ as input. The numerical values change by $\sim 10\%$
  (from $1.1$ to $1.12\ee{-10}$~m/s$^2$), which is within the $8\%$
  discrepancy already noted in W3i6.
\item \textbf{P14 $\to$ P5}: Timing predictions in Sector~A use
  $\psigrav = GM/c^2r$, unaffected by the MOND formula or the
  $\alpha^{57}$ chain.
\item \textbf{P15 $\to$ P2/P11}: Channel~B (now highest priority)
  requires P2 toroidal cavity + P11 SRF coating technology.
  \textbf{Recommend accelerating P2/P11 R\&D.}
\item \textbf{Cosmo $\to$ all}: T4 tentative resolution
  (Finding~\ref{find:void-evolution}) reduces the most serious tension
  from $3$--$5\times$ to $1$--$2\times$, significantly strengthening
  the overall DFD case.
\item \textbf{P14 $\to$ MatSci}: No changes. YBCO FOM = 965,
  Nb$_3$Sn FOM = 92, all material recommendations stand.
\end{itemize}

%=============================================================================
\section{Conclusion}
\label{sec:conclusion}
%=============================================================================

Wave~3 Iteration~7 advances the canonical DFD on three fronts while
identifying two significant errors in W3i6.

\textbf{Advances:}
\begin{enumerate}[nosep]
\item T4 Cold Spot over-prediction reduced from $3$--$5\times$ (W3i6)
  to $1$--$2\times$ (W3i7) by the discovery of the void-evolution
  partial cancellation effect. T4 is tentatively resolved.
\item $\alpha^{57}$ confirmed unaffected by the two-component
  decomposition (UV/IR separation argument).
\item $\mufd(x) = x/(1+x)$ derivation from $S^3$ composition law
  confirmed compatible with the two-component framework; $\kappa$
  rendered operationally irrelevant.
\item Detailed ESS measurement protocol provided with explicit
  5-sigma criteria.
\end{enumerate}

\textbf{Corrections:}
\begin{enumerate}[nosep]
\item The W3i6 MOND formula $\aM = c\Hzero\sqrt{\Omega_\Lambda}$ is
  numerically incorrect ($5.4\ee{-10}$ vs.\ claimed $1.1\ee{-10}$).
  Revert to $\aM = 2\sqrt{\alpha}\,c\Hzero$.
\item The W3i6 Channel~A sensitivity $\lamdiff_{\rm min} = 1.14\ee{-14}$
  is incorrect. Corrected value: $\sim 10^{-7}$. Channel~A cannot
  improve the current SQMS bound.
\end{enumerate}

The canonical DFD now has: 4 axioms, effectively 1 operational free
parameter ($\lambda$), 13+ testable predictions, all 5 tensions
resolved or tentatively resolved, and a clear experimental programme
with Channel~B as the highest-priority bound-improving experiment.

\end{document}
